\documentclass[prep,both]{debate}
\motion{语言的边界是 / 不是人类的边界}
\meta{赛制}{招新赛 3v3}
\meta{生成日期}{2026-09-30}
\begin{document}
\debatetitle

\begin{summary}
\begin{fields}
\field{持方优劣评估}{整体略偏反方，\textbf{正方较劣势}。劣势主要来自论证义务和评委直觉：“是”字要求两条边界基本重合，而近十年的神经科学（失语研究、Fedorenko 等 2024 年的 Nature 综述）站在“思考不以语言为前提”一边；“言不尽意”“只可意会不可言传”又是中文听众的常识。正方的核心打法是把战场收窄到两处：人超出动物的那一截（精确的数、理解别人想错了），以及人类一代接一代的积累；反复讲清“一个人摸到”和“人类走到”的区别，把反方的失语者、默会知识都归到“摸到”。反方要防的是滑成“语言不重要”，以及被“脚手架”一击打回。}
\thesis{正方}{人比别的动物多走出去的每一步，都是语言铺的路。语言到不了的地方，一个人也许摸得到，人类走不到。}
\thesis{反方}{语言记下的是人走过的路。人总是先想到、先做到，词才跟上来；所以人能走到的地方，比语言能说到的地方远。}
\end{fields}
\end{summary}

\section{一、赛制要点}
\begin{flowtable}
\block{A 立论质询}
\flow{1 正方一辩立论}{180 秒}{正一}{定义、中立判准、两个论点讲清；抛出全场问题}
\flow{2 反方二辩质询正方一辩}{90 秒 双边}{反二 问，正一 答}{反方：拿到“失语的人还能思考”的承认；正方：先答是否，守住“摸到与走到”}
\flow{3 反方一辩立论}{180 秒}{反一}{开头用一句话接住质询所得，再立两个论点}
\flow{4 正方二辩质询反方一辩}{90 秒 双边}{正二 问，反一 答}{正方：拿到“失语者是先有语言的”承认；反方：守住“思考不是用语言做的”}
\block{B 三辩对辩与二辩申论}
\flow{5 双方三辩对辩}{各 90 秒}{正三 与 反三}{开第一个战场：思考。尼加拉瓜手语者对失语者}
\flow{6 反方二辩申论}{90 秒}{反二}{拆正方立论主干 + 推进“想在说前面”（爱因斯坦的信）}
\flow{7 正方三辩质询反方二辩}{90 秒 双边}{正三 问，反二 答}{追问爱因斯坦的“第二阶段”：没找到词，人类知不知道}
\flow{8 正方二辩申论}{90 秒}{正二}{回应反二申论 + 拆反方立论主干 + 推进“没有词，走不到”（家庭手势者）}
\flow{9 反方三辩质询正方二辩}{90 秒 双边}{反三 问，正二 答}{追问“那门手语是谁造的”}
\block{C 一辩对辩与三辩申论}
\flow{10 双方一辩对辩}{各 90 秒}{正一 与 反一}{收拢前两段：“摸到”与“走到”}
\flow{11 正方三辩申论}{90 秒}{正三}{处理反方最强点 + 推进“说不出就接不上”（性骚扰一词）}
\flow{12 反方一辩质询正方三辩}{90 秒 双边}{反一 问，正三 答}{为反方结辩取材}
\flow{13 反方三辩申论}{90 秒}{反三}{回应正三申论 + 处理正方最强点 + 推进“知道的比说的多”（俄罗斯蓝）}
\flow{14 正方一辩质询反方三辩}{90 秒 双边}{正一 问，反三 答}{全场最后一次质询，为正方结辩取材}
\block{D 二辩对辩与结辩}
\flow{15 双方二辩对辩}{各 90 秒}{正二 与 反二}{结辩前最后交锋：守住中立判准的“稳定到达”}
\flow{16 反方一辩结辩}{90 秒}{反一}{归纳分歧、处理最强点、封路、价值升华}
\flow{17 正方一辩结辩}{90 秒}{正一}{回应反一结辩、约 40 字临场、价值升华}
\end{flowtable}
质询 90 秒双边计时，频繁刻意打断要扣分：两条链共四到七问，每问一句话答得完，对方答非所问先复述一次再问。没有驳论、小结和自由辩，对方立论要在两篇申论里拆，两篇分工写死（见第十节口径表）。一辩稿算全队作品，一辩的个人分来自被质询、一辩对辩、质询对方三辩和结辩，这几处的材料要给足。

\section{二、辩题性质与争议焦点}
\begin{fields}
\claim[辩题性质]{事实判断，带哲学性。“是 / 不是”判断两条边界是否重合。题目化用维特根斯坦《逻辑哲学论》5.6，但把“我的”换成了“人类的”，所以不能只靠解释维特根斯坦来赢。}
\field{正方倾向把题目解释成}{一个关于条件的命题：人之为人的那一截能力，以语言为条件；语言走到哪，人才走到哪。}
\field{反方倾向解释成}{一个关于存在的命题：只要在语言之外，人能稳定地思考、做到并守住重要的东西，两条边界就不重合。}
\field{争议焦点}{\sub{思考}{语言坏了、还没学会时，人还能不能想？失语者、婴儿、爱因斯坦，对尼加拉瓜手语者、家庭手势者。}\sub{积累}{说不出的东西能不能进入“人类”？轮扁、石器实验、棘轮效应，对默会知识、示范教学。}\sub{方向}{边界往外扩时，是人先到词后到，还是词到了人才到？性骚扰一词、爱因斯坦的信、海伦·凯勒的“无世界”。}}
\end{fields}

\section{三、关键词定义}
\subsection{3.1 语言}
\begin{fields}
\field{调研}{\sub{词典义}{人类用来表达思想、进行交际的工具，由语音、词汇和语法构成（汉典“语言”条）。}\sub{学术义}{语言学与认知科学说的是自然语言，口语、书面语、手语都算；Fedorenko 等 2024 所说的大脑“语言网络”对三者同样反应，Pyers \& Senghas 2009 的研究对象就是尼加拉瓜手语。}\sub{哲学义}{《逻辑哲学论》里的“语言”是命题的总体，是逻辑意义上能有意义地说出的东西；日常说“音乐的语言”“身体语言”是比喻。}}
\begin{procard}{正方有利定义}\definition{语言，是人类一切用符号表达、传递意义的系统，包括自然语言、数学、乐谱、图像和手势。}反方举的数学、音乐、手艺示范都被收进“语言”，边界自然重合。\sub{有利之处}{反例几乎全部失效。}\sub{反方如何反驳它}{定义里藏了结论：什么表达都叫语言，题目就成了“人能表达的边界是人能表达的边界”；词典要求有词汇和语法；音乐叫语言只是比喻。评委会认为这是定义杀。}\end{procard}
\begin{concard}{反方有利定义}\definition{语言，是嘴上说出来、纸上写下来的话。}手语、手势、数学、音乐、手艺都落在语言之外，反例随手可得。\sub{有利之处}{举例成本极低。}\sub{正方如何反驳它}{把手语赶出语言，就等于说聋人没有语言，这是语言学早已否定的看法；辩题问的是语言本身能到哪里，不是嘴能到哪里。}\end{concard}
\begin{neutralcard}{中立定义}语言，是有词、有语法、用来表达和交流意思的符号系统。口语、文字、手语都是；数学记号这类靠语言定义、从语言里长出来的书写系统算语言的延伸；音乐、绘画、表情、手上的功夫叫“语言”，只是比喻。\sub{合理性}{贴近词典（有词汇、有语法）；贴近语言学共识（手语是语言）；双方都有得打：正方拿得到数词和心理词，反方拿得到音乐、手艺和失语者的思考。}\end{neutralcard}
\end{fields}

\subsection{3.2 边界（重点是“人类的边界”）}
\begin{fields}
\field{调研}{\sub{词典义}{边界：国家或地区之间的分界线；界限：限度、止境，也指事物之间的分界（汉典）。}\sub{两种读法}{“止境”读法问人能走多远；“分界”读法问人和非人的线画在哪。“语言的边界”双方都读成语言能说清的范围到哪里为止，争的是“人类的边界”。}}
\begin{procard}{正方有利定义}\definition{人类的边界，是人类作为一个整体能共同认识、积累和传承的范围；也可以读成人与其他动物的分界线。}\sub{有利之处}{个人说不出的感受不算“人类”的边界；动物也有的能力不算人之为人的那条边。}\sub{反方如何反驳它}{人类由一个个人组成，一个人想到、做到的，人类就到了；“分界线”读法把题目偷换成“语言是不是人的本质”，那是另一道题。}\end{procard}
\begin{concard}{反方有利定义}\definition{人类的边界，是人能感受、体验、做到的一切的极限，包括感觉、情绪和身体技能。}\sub{有利之处}{疼、爱、骑车的平衡、认脸，人人都有又说不清，边界显然不重合。}\sub{正方如何反驳它}{狗也会疼，也认得主人。照这个定义，人的边界和狗差不多，题目里“人类”两个字被抹掉了；感受到不等于走到。}\end{concard}
\begin{neutralcard}{中立定义}人类的边界，是人能想到、认识到、做到，并且能传给别人的范围的极限。既包括一个人的思考和本事，也包括人类共同的积累。\sub{合理性}{贴近“止境”义；个人与集体都在里面：反方可以用个人的思考和手艺，正方可以用积累和传承。}\end{neutralcard}
\end{fields}

\subsection{3.3 是}
\begin{fields}
\begin{procard}{正方有利定义}\definition{“是”，就是划定：语言的边界决定人类的边界，不要求一点不差。}\sub{有利之处}{零星反例不影响结论。}\sub{反方如何反驳它}{题目写的是“是”，不是“影响”；把“是”降成“影响”，就只剩“语言很重要”，那不需要辩。}\end{procard}
\begin{concard}{反方有利定义}\definition{“是”，就是处处重合：只要有一处人能到而语言到不了，就不是。}\sub{有利之处}{举一个说不清的感受就够。}\sub{正方如何反驳它}{拿一个说不清的瞬间取消整道题；任何边界都有毛边，国界上也有飞过去的鸟。}\end{concard}
\begin{neutralcard}{中立定义}“是”，指两条边界基本重合：语言到不了的地方，人类在重要的方面也到不了。不要求逐点重合，但反例必须稳定、成体系、对人重要，不能是一闪而过的说不清的感受。\sub{合理性}{既不让正方把“是”降成“影响”，也不让反方拿一个瞬间取消题目。}\end{neutralcard}
\end{fields}

\subsection{3.4 隐含要素}
\begin{fields}
\field{主体}{人类（物种）以及其中的每个人，两个层次都算。这是正方“摸到与走到”、反方“一个人先到，人类就到了”的交锋点。}
\field{范围}{普遍命题，不限时代和地域。}
\field{时间尺度}{两条边界都在移动：新词、新知识、新技术。比较的是它们是否同步、谁带着谁走。}
\field{比较对象}{语言能到的范围，对人能到的范围；不是“有语言的人和没语言的人谁过得好”。}
\end{fields}

\begin{warnbox}{3.5 定义杀陷阱：双方都要背}
\begin{fields}
\begin{procard}{正方的定义杀陷阱}把语言扩成“一切有意义的表达”，或者把人类的边界缩成“人与动物的分界”再说“语言是人的本质”。评委会觉得正方在用定义赢，被点破一次就很难挽回。正方的数字、乐谱只在“数学记号是语言的延伸”这一句里用，音乐本身承认是比喻。\end{procard}
\begin{concard}{反方的定义杀陷阱}把语言缩成“嘴上说的话”，把手语、文字踢出去；或者把“是”抬成“逐点重合”，然后举一个“我说不出有多爱我妈”就收工。评委会认为反方在回避题目的分量。反方的例子一律用稳定、成体系的能力：失语者的算术、作曲家的作品、手艺人的手。\end{concard}
\end{fields}
\end{warnbox}

\section{四、判准与论证义务}
\begin{fields}
\claim[判准是什么]{在这道题里，评委看的是：语言到不了的地方，人类还能不能稳定地到达。}
\begin{procard}{正方有利判准}\definition{看人类共同拥有的东西：凡是能被检验、积累、传给下一代的成果，是不是都要经过语言。}\sub{有利之处}{个人一闪而过的感受、动物也有的能力都不算；失语者、婴儿的例子被挡在外面。}\sub{反方如何反驳它}{把“个人”开除出“人类”；而且手艺、音乐靠示范就能传，不经过语言照样守了几千年。}\end{procard}
\begin{concard}{反方有利判准}\definition{看有没有语言之外的人类能力：只要找到人不用语言仍然做得到的重要事情，边界就不重合。}\sub{有利之处}{失语者算数、作曲，一条就够。}\sub{正方如何反驳它}{只证明了“有”，没证明这些能力不是语言先搭起来的：失语者是先学会语言、后来才失去的；婴儿会的一加一，猴子也会。}\end{concard}
\begin{neutralcard}{中立判准}看语言到不了的地方，人类还能不能稳定地到达。\sub{合理性}{直接对应“边界是否重合”；不偏向个人或集体；把“偶尔说不清”和“稳定地到达”分开，双方都要证明实质的东西。}\sub{正方需要证明}{人超出本能的那一截（精确的思考、理解别人、积累知识）以语言为条件；语言之外的东西，要么到不了，要么到了也留不住、接不上。只证明“语言很有用”不够。}\sub{反方需要证明}{在语言之外，人能稳定地思考、认识、创造，而且这些东西对人重要、留得住（至少能被一再做到，最好能传下去）。只证明“语言不完美、有说不清的感受”不够。}\end{neutralcard}
\end{fields}

\prosection{五、正方论点}
\prosubsection{论点一：没有词，走不到}
\begin{fields}
\claim{人超出动物的那一截思考，精确的数、别人心里想错了什么，是跟着语言一起到的。语言走到哪，人才走到哪。}
\begin{mechanism}\step 人天生只有和动物共享的底子：三以内的数、大概多少、眼前的东西。\step 词把这些底子固定下来、组合起来，带过时间和空间：“十七”能被记住，“他以为”能被想。\step 一辈子缺这个词的人，哪怕身在会数数的社会、有二十多年社交经验，也走不到这一步。\end{mechanism}
\begin{evidence}{学理}[已核实]维果茨基《思维与言语》（1934）第 7 章：思想不是在词里表达出来的，而是在词里完成的。\sub{核心主张}{言语不是给想好的思想穿衣服，思想在变成言语时被重组、被完成。}\sub{支撑}{机制第二步：词不只记录思想，还完成思想。}\sub{边界}{维果茨基同一章也说思想和词“不是一个模子刻出来的”，承认两者根源不同。正方只主张“高级的那一截”，不主张“没有语言就没有思考”。}\sub{出处}{Vygotsky, Thinking and Speech, ch.7, Minick 英译：“Thought is not expressed but completed in the word.”}\sub{口头版}{维果茨基说，思想是在词里完成的。}\end{evidence}
\begin{evidence}{数据}[已核实]尼加拉瓜手语第一批使用者 8 人（2001 年平均 26.8 岁），低语言的错误信念任务 4 题平均只答对 0.5 题；第二批 10 人（平均 17.5 岁）平均答对 3.8 题。两年后，学到心理词汇的第一批明显进步。作者结论：语言是必要前提，“连二十五年的社会经验也替代不了”。\sub{来源}{Pyers, J. E. \& Senghas, A. (2009). Language promotes false-belief understanding: Evidence from learners of a new sign language. Psychological Science, 20(7), 805–812. doi:10.1111/j.1467-9280.2009.02377.x}\sub{方法}{两年追踪（2001、2003），聋校学生与校友；图片补全式的低语言错误信念任务；同时记录心理状态手语词的使用。}\sub{局限}{样本很小（8 人对 10 人）；追踪相关而非随机实验；第一批两年后仍只答对一半左右。}\end{evidence}
\begin{evidence}{数据}[已核实]4 名尼加拉瓜成年家庭手势者（20–29 岁，没有任何可用的语言输入，但有工作、会挣钱、认得钱的面值）：目标是 1 到 3 个时 100\% 准确；超过 3 个，数量接近但精确答对的不到一半。\sub{来源}{Spaepen, E., Coppola, M., Spelke, E. S., Carey, S. E. \& Goldin-Meadow, S. (2011). Number without a language model. PNAS, 108(8), 3163–3168. doi:10.1073/pnas.1015975108}\sub{方法}{对成年家庭手势者做数量交流与集合配对任务，与有语言的对照组比较。}\sub{局限}{只有 4 人；他们估得接近，说明缺的是精确，不是大概，反方会用。}\sub{用在}{正二申论“推进论点一”。}\end{evidence}
\begin{evidence}{案例}[已核实]尼加拉瓜第一批手语者：成年人，能工作、会交朋友，却答不出“那个人以为球在哪个盒子里”；学会“想”“知道”一类的手语词以后，这道题开始答对。\sub{代表性}{特殊人群；正因为他们只缺语言、不缺社会经验，才把语言的作用单独拎了出来。}\sub{出处}{同上 Pyers \& Senghas 2009。}\end{evidence}
\closing{精确的数、理解别人想错了什么，是人超出本能的那一截。这一截语言到不了，人就到不了，所以在人之为人的地方，两条边界重合。}
\begin{clash}\pair{失语的人什么都说不出，照样算数、下棋、作曲。}{楼盖好了拆掉脚手架，楼还在，不能说明楼不靠脚手架盖。他们都是先学会语言、几十年后才失去的；Fedorenko \& Varley 2016 自己写了这个可能：就算语言在发育中是关键，成熟的大脑里也可能不再需要。一辈子缺语言的人，才是检验，他们走不到。}\pair{尼加拉瓜那门手语，是一群没有语言的孩子自己造的，人走在语言前面。}{在造出“想”“知道”之前，他们走不到那一步。边界是一起挪的，挪的办法只有一个：造词。}\end{clash}
\end{fields}

\prosubsection{论点二：说不出就接不上}
\begin{fields}
\claim{人类的边界，是一代接一代往前走的边界。说不出的东西能被模仿、被守住，却接不上力，跟着一个人活，也跟着一个人走。}
\begin{mechanism}\step 人类靠一代接着一代往前走（棘轮效应）。\step 接力要传得准，最准、能被检验和纠正的是语言。\step 说不出的只能模仿：走样、停滞、随人而去。\end{mechanism}
\begin{evidence}{学理}[已核实]棘轮效应（Tennie, Call \& Tomasello, 2009）\sub{核心主张}{人类文化会随时间累积修改，黑猩猩的传统只在已有能力范围内打转。}\sub{支撑}{机制第一步。}\sub{边界}{作者把原因归于“重过程的社会学习”和主动教学、从众规范，不只是语言。正方用它说明“接力”是人类的特征，语言的作用靠石器实验补。}\sub{出处}{Tennie, C., Call, J. \& Tomasello, M. (2009). Ratcheting up the ratchet. Phil. Trans. R. Soc. B, 364, 2405–2415.}\sub{口头版}{人类的本事是一代接一代往上攒。}\end{evidence}
\begin{evidence}{数据}[已核实]184 名成年人排成传递链学打奥杜威石器。只看师傅打、不能交流，每一下敲出可用石片的概率没有提高；只许比划不许出声，翻一倍；允许说话，翻四倍（都相对只看石片的基线）。作者还指出奥杜威技术停滞了七十多万年（原文 more than 700,000 years），推测与传得不准有关。\sub{来源}{Morgan, T. J. H. et al. (2015). Experimental evidence for the co-evolution of hominin tool-making teaching and language. Nature Communications, 6, 6029. doi:10.1038/ncomms7029}\sub{方法}{实验室传递链：五种条件（逆向工程、模仿、基础教学、手势教学、言语教学），每种六条链，测量石片数量、质量和每下成功率。}\sub{局限}{现代人模拟古人；“停滞与传得不准有关”是作者推测；比划也有提高，反方会用。}\end{evidence}
\begin{evidence}{数据}[已核实]连反方最常引用的那篇 Nature 综述也写：语言是传递文化知识的强大工具。\sub{来源}{Fedorenko, E., Piantadosi, S. T. \& Gibson, E. A. F. (2024). Language is primarily a tool for communication rather than thought. Nature, 630, 575–586. doi:10.1038/s41586-024-07522-w}\sub{局限}{这篇的主结论是“语言不是复杂思考的前提”，正方只能引它的“传递”一句，不能说它支持正方。}\end{evidence}
\begin{evidence}{案例}[已核实]轮扁斫轮（《庄子·天道》）：七十岁的老木匠说，做轮子“得之于手而应于心，口不能言，有数存焉于其间。臣不能以喻臣之子，臣之子亦不能受之于臣……古之人与其不可传也死矣”。\sub{代表性}{寓言；它正是反方最爱举的“不可言传”，而它自己写出了代价：说不出的，儿子学不走，跟着人一起死了。}\sub{出处}{《庄子·天道》，维基文库本。}\end{evidence}
\begin{evidence}{案例}[已核实]“性骚扰”这个词：1975 年，康奈尔大学物理系的行政人员 Carmita Wood 被上司骚扰多年后辞职；Lin Farley 等人和她开会讨论，最后定下了“sexual harassment”这个说法。Farley 说：“我觉得我们需要给这个现象一个名字，我们都得在说同一件事。”这个词随后进入全国词汇，也进入了法律。\sub{代表性}{一个名字把千万人各自的遭遇变成人类的共同知识。}\sub{出处}{Blakemore, E. (2018-01-08). Until 1975, “Sexual Harassment” Was the Menace With No Name. History.com；学理见 Fricker 2007 诠释不正义。}\sub{用在}{正三申论“推进论点二”。}\end{evidence}
\closing{人类的边界是能接着往前走的边界。能接上的，都要经过语言；接不上的，留在一个人那里，他走了就没了。所以在人类这个层面，两条边界重合。}
\begin{clash}\pair{奥杜威石器没有像样的语言，照样传了七十多万年。}{传了，也停了七十多万年。守住和往前走是两回事，人类的边界是往前走出来的。}\pair{比划组也翻了一倍，不说话也教得会。}{一倍对四倍。比划能帮人守住，说话才让人往前走得动。}\end{clash}
\end{fields}

\subsection{（被淘汰的正方候选论点，交给反驳库与后场备用）}
\begin{enumerate}[leftmargin=1.8em]
\item 社会实在只活在语言里：法律、货币、权利由“X 算作 Y”的宣告造出来（Searle 2010）。只有哲学论证，没有数据，备用。
\item 语言是人与非人的分界：定义杀风险，反方会说正方偷换题目。
\item 语言塑造感知（俄罗斯蓝）：效应只是快了一百多毫秒，英语者照样九成六分得清，会被反用。
\item 界限只能在语言里划（《逻辑哲学论》序言）：哲学味太重，留作回应“维特根斯坦承认有不可说的”。
\item 语言让成人整合信息（Hermer-Vazquez 等 1999，边跟读边找方向的成人表现得像幼儿和老鼠）：Ratliff \& Newcombe 2008 三个实验没有复现，结论是“语言不是必需的”。不上场。
\end{enumerate}

\consection{六、反方论点}
\consubsection{论点一：想在说前面}
\begin{fields}
\claim{思考不是用语言做的。语言坏了，人照样算数、推理、作曲；语言还没到，人已经想到了。}
\begin{mechanism}\step 思考和语言用的是大脑里两套不同的系统。\step 所以语言系统坏了，算数、推理、作曲都还在。\step 所以人能先想到语言还说不出的东西，再回头去找词。\end{mechanism}
\begin{evidence}{学理}[已核实]语言主要是交流的工具，不是思考的工具（Fedorenko, Piantadosi \& Gibson, 2024, Nature）\sub{核心主张}{大脑里的语言网络与负责数学、逻辑、推断他人想法的系统双向分离；语言“不像是复杂思考、包括符号思考的前提”。}\sub{支撑}{机制第一、二步。}\sub{边界}{同文承认语言剥夺的孩子“有些推理确实延迟”，也承认语言是传递文化知识的强大工具。反方只用“不是前提”，不说“语言对思考没用”。}\sub{出处}{Nature 630, 575–586. doi:10.1038/s41586-024-07522-w}\sub{口头版}{语言是说给别人听的工具，不是用来想的材料。}\end{evidence}
\begin{evidence}{数据}[已核实]三名左脑大面积损伤、重度语法失语的男性（摘要未说明病因），所有基本运算完好，能解带括号、要按结构运算的算式。全面失语者几乎完全不能理解和说出语言，却仍能加减、解逻辑题、推断别人的想法、欣赏音乐。\sub{来源}{Varley, R. A. et al. (2005). Agrammatic but numerate. PNAS, 102(9), 3519–3524. doi:10.1073/pnas.0407470102；Fedorenko, E. \& Varley, R. (2016). Language and thought are not the same thing. Annals of the NYAS, 1369, 132–153. doi:10.1111/nyas.13046}\sub{方法}{神经心理学病例研究（三例）；综述病例研究与脑成像研究。}\sub{局限}{病例少；病人发病前有语言，正方会打“脚手架”。}\end{evidence}
\begin{evidence}{数据}[已核实]爱因斯坦 1945 年回答数学家阿达马的问卷：写出来或说出来的词语，在他的思考机制里似乎不起任何作用；思考的元素是可以自主再现、组合的符号和图像，约定的词语或其他符号，要到第二阶段才费力去找。\sub{来源}{Hadamard, J. (1945). The Psychology of Invention in the Mathematical Field. Princeton University Press（Dover 1954 重印），附录二，第 142–143 页。}\sub{局限}{一个人的内省，个例，只用来展示机制。}\sub{用在}{反二申论“推进论点一”。}\end{evidence}
\begin{evidence}{案例}[已核实]苏联作曲家舍巴林中风后严重失语，继续作曲，作品被评为与病前水准相当。\sub{代表性}{个案，展示“语言坏了、创造还在”的机制。}\sub{出处}{Luria, Tsvetkova \& Futer (1965). Aphasia in a composer (V. G. Shebalin). J. Neurol. Sci. 2, 288–292；转述见 Fedorenko \& Varley 2016。第五交响曲与肖斯塔科维奇的评价只见百科，待核实，不上场。}\end{evidence}
\closing{思考不靠语言做，语言坏了思想还在，思想还能跑在词前面。人能想到的地方，语言框不住，所以两条边界不重合。}
\begin{clash}\pair{失语者是先学会语言、后来才失去的，楼盖好了才拆脚手架。}{Fedorenko 2024 同一篇写：孩子身上语言和思考的分离早就存在；语言剥夺的聋孩子照样学会数学、做因果推理。只有延迟，没有边界。}\pair{爱因斯坦要是没找到词，相对论就只在他脑子里。}{他先到了，这就说明人的边界在词前面。词是后来去接他的，不是先去划他的。}\end{clash}
\end{fields}

\consubsection{论点二：知道的比说的多}
\begin{fields}
\claim{人最高处的本事，棋感、手上的分寸、医生的直觉，说不清，却真实存在，靠做和示范一代代传下去。说不出，不等于到不了。}
\begin{mechanism}\step 很多知识存在手里、眼里，用得出，说不出。\step 它们靠示范、手把手、一起做来传。\step 所以人能到、能守住的地方，比语言能说到的地方大。\end{mechanism}
\begin{evidence}{学理}[有把握]默会知识（波兰尼，1966，《默会的维度》）\sub{核心主张}{“我们知道的比我们能说出来的多。”我们能在一千张、一百万张脸里认出一张，却说不出是怎么认出来的。}\sub{支撑}{机制第一步。}\sub{边界}{默会知识常和说得出的知识交织：师傅也会说“腕子松一点”。反方只主张“有说不出的那一截”，不主张全部说不出。}\sub{出处}{Polanyi, M. (1966). The Tacit Dimension. Chicago: University of Chicago Press, ch.1。本次只核到二手引用，赛前请查原书第 4 页。}\sub{口头版}{我们知道的，比我们说得出的多。}\end{evidence}
\begin{evidence}{数据}[已核实]同一个石器实验里，只许比划、不许出声的教学组，每一下敲出可用石片的概率翻了一倍。\sub{来源}{Morgan, T. J. H. et al. (2015). Nature Communications, 6, 6029.}\sub{局限}{说话组翻四倍，正方必用；反方的口径是“快是工具好用，不是边界”。}\end{evidence}
\begin{evidence}{数据}[已核实]俄语把浅蓝和深蓝分成两个词，英语只有一个“蓝”。让 26 个俄语者、24 个英语者分辨同一组蓝色：准确率英语者 96.5\%、俄语者 95.7\%；俄语者跨词的两种蓝分得更快，类别优势 124 毫秒，一做言语干扰优势就没了。\sub{来源}{Winawer, J. et al. (2007). Russian blues reveal effects of language on color discrimination. PNAS, 104(19), 7780–7785. doi:10.1073/pnas.0701644104}\sub{方法}{限时辨色实验，加言语干扰与空间干扰对照。}\sub{局限}{双刃：它也证明语言会实时影响感知速度。反方只用“英国人没有这个词也分得清”。}\sub{用在}{反三申论“推进论点二”。}\end{evidence}
\begin{evidence}{案例}[有把握]认脸：你能在车站一千张脸里一眼认出你妈，却说不出她的鼻子多高、眼睛多宽。\sub{代表性}{人人都有，评委能马上想到自己。}\sub{出处}{波兰尼 1966 的原例。}\end{evidence}
\closing{人能做到、能守住的本事，有一大截语言说不出。说不出却到得了，所以两条边界不重合。}
\begin{clash}\pair{轮扁的儿子学不走，“古之人与其不可传也死矣”。}{轮扁七十岁还在斫轮，人到得了。他说的“不能以喻”，是用话教不会；手艺本来就不靠话教。}\pair{俄国人快了一百多毫秒，语言在划线。}{快一点慢一点是语言在调速度；分得清分不清才是边界。英国人没有这个词，照样九成六。}\end{clash}
\end{fields}

\subsection{（被淘汰的反方候选论点，交给反驳库与后场备用）}
\begin{enumerate}[leftmargin=1.8em]
\item 人先到词后到：并入论点一，作反二申论的推进材料。
\item 婴儿与动物：五个月的婴儿会算一加一（Wynn 1992）。并入论点一备用；正方会说“猴子也会”。
\item 语言相对论只剩弱版本：沃尔夫说霍皮语里没有指时间的词，Malotki 1983 用六百多页的实例反驳。只打对方，不立我方，进反驳库。
\item 维特根斯坦自己承认有不可说的：6.522、致 Ficker 信。只拆对方引用，进质询与反驳库。
\item 艺术与音乐：并入论点二。
\end{enumerate}

\prosection{七、正方反驳库（针对反方全部可能论点）}
\begin{rebuttals}
\rebuttal{1}{想在说前面：失语的人照样算数、推理、作曲。}{楼盖好了拆掉脚手架楼还在，不说明楼不靠脚手架盖。}{失语者都是先学会语言、几十年后才失去的。Fedorenko \& Varley 2016 自己写：就算语言在发育中是关键，成熟的大脑里也可能不再需要。真正的检验是一辈子没有语言的人：尼加拉瓜第一批手语者二十多年社交也答不对错误信念题。}{对方：“Fedorenko 2024 说语言剥夺的聋孩子也能学数学。” → 同一段写“有些推理确实延迟”；而且他们是在后来接触了数字和符号之后学的。}{最终}
\rebuttal{2}{知道的比说的多：棋感、手艺、医生的直觉说不清。}{一个人知道的比说的多，人类知道的不比说的多。}{说不出的那一点，每一代都要从零摸。轮扁自己说“臣不能以喻臣之子……古之人与其不可传也死矣”。石器实验里只模仿的那组，没有提高。}{对方：“手艺传了几千年。” → 传下来的是能说的那部分和一双双重新摸的手；说不出的那点每代重来，这就是边界。}{最终}
\rebuttal{3}{爱因斯坦先有图像，后找词。}{他先摸到，人类要等他找到词才走到。}{信里他自己写，词语要到第二阶段才费力去找，为的是能传达给别人。没找到，相对论就只在他一个人脑子里。}{对方：“那人类的边界就在他脑子里。” → 在他脑子里的是他的边界，写成论文以后才是人类的。}{最终}
\rebuttal{4}{俄罗斯蓝：英国人没有两个词照样分得清。}{我方不靠感知立论，感知是人和动物共有的底子。}{对方自己的数据显示俄国人更快，言语干扰一来优势就没了，语言在实时参与。但我方讲的是超出底子的那一截。}{对方：“那你们承认语言划不住感知。” → 承认。看见颜色不是人类的边界，给颜色分类、讲出光谱才是。}{最终}
\rebuttal{5}{婴儿五个月就会算一加一。}{猴子身上也测到过（Hauser 等 1996），那是动物的边界；婴儿要等学会数词才到得了十七。}{Wynn 1992 是三以内的小数；Spaepen 2011 的家庭手势者过了三就不准。}{对方：“那是天生的思考。” → 对，天生的底子，人和动物共有；题目问的是人类的边界。}{淘汰}
\rebuttal{6}{维特根斯坦自己说有不可说的（6.522，致 Ficker 信）。}{他让我们对不可说的保持沉默；沉默的东西，一个人可以活出来，人类没法接着往前推。}{他在序言里说，界限只能在语言中划出。我方不靠他立论，用的是尼加拉瓜和石器。}{对方：“那 5.6 你们不用了？” → 5.6 说的是“我的”，我方论证的是人类。}{淘汰}
\rebuttal{7}{5.6 说的是“我的语言”“我的世界”，不是人类。}{同意，所以我方论证的是人类，用的是实验，不是一句名言。}{辩题把“我的”换成了“人类的”，这一换，问题就从一个人摸到哪里，变成人类走到哪里。}{对方：“那辩题的出处就不支持你。” → 辩题问的是事实，不是出处。}{淘汰}
\rebuttal{8}{言不尽意、只可意会不可言传。}{说不尽的那一半，就留在了一个人那里，这正是我方说的边界。}{《系辞》下一句是“圣人立象以尽意”，古人想的是怎么把意说尽，不是守着说不尽。}{对方：“象不是语言。” → 象要靠辞来系，所以下一句就是“系辞焉以尽其言”。}{常见}
\rebuttal{9}{音乐、艺术超出语言。}{一首曲子能传三百年，靠的是乐谱和乐理；演奏者手上那一点，每代重练。}{我方承认音乐不是语言；我方说的是，能被接着往前推的，是被记下来、讲出来的那部分。}{对方：“乐谱不是语言。” → 乐谱是靠语言定义的记号，按中立定义算语言的延伸。}{常见}
\rebuttal{10}{尼加拉瓜孩子自己造了语言，人走在语言前面。}{他们在造出“想”“知道”之前走不到那一步。}{边界一起挪，挪的办法只有造词。反方举的恰恰是“有了词才走到”的例子。}{对方：“造词的是人。” → 人造词，词带人走，这是一起挪；没有哪一步是人不带词先走过去的。}{常见}
\rebuttal{11}{动物没有语言也会思考（乌鸦、猩猩）。}{思考有语言之外的底子，那是动物和人共有的。}{人类比乌鸦多出来的那一截，是语言铺的：精确的数、别人心里想错了什么、一代接一代的积累。}{对方：“那你承认语言外面有思考。” → 承认有底子，不承认有人类的边界在那里。}{常见}
\rebuttal{12}{爱和疼说不清，人的世界远大于语言。}{狗也会疼、也会依恋，感受到不等于走到。}{人类为疼造了医学的词，为爱写了几千年的诗，这正是语言在往外推。}{对方：“那说不清的那部分呢？” → 留在你心里，也就只在你心里。}{常见}
\end{rebuttals}

\consection{八、反方反驳库（针对正方全部可能论点）}
\begin{rebuttals}
\rebuttal{1}{没有词，走不到：尼加拉瓜第一批手语者答不对错误信念题。}{这证明学词能帮人学会，不证明语言到不了的地方人到不了；造出这门语言的，正是一群没有语言的孩子。}{第一批只有 8 人。Fedorenko 2024 写，语言剥夺的人“有些推理确实延迟”，延迟不是边界。第二批学会就答对了，词是人造的。}{对方：“二十五年的社交也替代不了。” → 替代不了的是一套心理词汇的教学；人能自己造出这套词，说明人的边界在语言外面。}{最终}
\rebuttal{2}{说不出就接不上：石器实验说话翻四倍，奥杜威停滞七十多万年。}{七十万年，是没有像样的语言照样传了七十万年；说话让人走得快，不等于只有说话才能走。}{比划组翻一倍，不说话也教得会。Fedorenko 2024 说语言“反映而不是产生”人类认知的精妙。}{对方：“说话是四倍。” → 快四倍是工具好用；汽车比走路快，公路不是人类的边界。}{最终}
\rebuttal{3}{轮扁：古之人与其不可传也死矣。}{轮扁七十岁还在斫轮，说明人到得了。}{他说的“不能以喻”，是用话教不会；手艺本来就不靠话教。}{对方：“儿子学不走。” → 学不走的是“用话教”这个办法，不是那门手艺。}{最终}
\rebuttal{4}{维果茨基：思想是在词里完成的。}{同一章他说思想和词“不是一个模子刻出来的”。}{维果茨基研究的是思想和言语在发育中怎样合流，他没说没有词就没有思想。}{对方：“完成也是一种依赖。” → 完成是说出来的那一刻；想到是在那之前。}{最终}
\rebuttal{5}{性骚扰：没有名字以前，受害者连自己都说不清。}{一九七五年以前她们就说得清：“他动手动脚，我不从，只好辞职。”缺的是一个名字，不是理解。}{这句原话出自造词的 Farley 本人。Fricker 2013 也承认另一种情形：本人完全清楚发生了什么，只是没法让机构听懂。}{对方：“没有名字就没有法律保护。” → 那是制度的边界，不是人的；名字是她们自己造的，人先到，词后到。}{淘汰}
\rebuttal{6}{社会实在活在语言里：货币、法律、国家。}{制度是人类的一部分，不是全部。}{我方不否认语言造了很多东西，问题是语言外面还有没有人。}{对方：“离开制度人还算人类吗？” → 失语的人离不开制度，也没有离开人类。}{淘汰}
\rebuttal{7}{语言塑造感知：俄国人分两种蓝更快。}{快了一百二十多毫秒是调速度，不是划边界。}{英国人没有两个词，准确率九成六，和俄国人一样。}{对方：“快慢也是边界。” → 边界是分得清和分不清。}{淘汰}
\rebuttal{8}{人和动物的分界就在语言。}{那是问语言是不是人的特征，不是问语言的边界是不是人的边界。}{会学舌的鹦鹉不是人，失语的舍巴林是。}{对方：“那人和动物靠什么分？” → 这道题不用我们回答这个。}{淘汰}
\rebuttal{9}{维特根斯坦：我的语言的界限意味着我的世界的界限（5.6）。}{两难：若是逻辑，5.641 明说哲学的“我”不是人；若是经验，失语者就是反例。}{6.522 他承认有不可说的东西，它显示自身；致 Ficker 的信说，没写出来的那部分才重要。}{对方：“他说的是逻辑形式。” → 那就和人类的边界无关，正方不能拿来用。}{常见}
\rebuttal{10}{海伦·凯勒：老师来之前，她活在一个“无世界”里。}{同一页她写，那时她会生气、会满足、会渴望，靠触觉记得这一切。}{她说的“无世界”是没有自我意识的世界，不是没有人；那段经验在语言之外存在过，后来才被她写出来。}{对方：“她是有了语言才写得出来。” → 写得出来，说明那段经验一直在，只是等语言去接。}{常见}
\rebuttal{11}{皮拉罕人没有数词，数量算不清。}{东西摆在眼前，他们一个不差地配对。}{研究者自己的结论：数词不改变底层的数量表征，只是帮人跨时间、跨空间记住。}{对方：“记不住就是边界。” → 记不住是工具没带，不是人到不了。}{常见}
\rebuttal{12}{界限只能在语言里划出（《逻辑哲学论》序言）。}{界限只能在语言里“划出来”，不等于只有语言里的才存在。}{6.522：不可说的东西存在，它显示自身。}{对方：“显示也要在语言里显示。” → 是在生活里显示，所以他说要对它沉默。}{常见}
\end{rebuttals}

\section{九、持方优劣评估}
\begin{assessment}
\assess{定义空间}{中立定义下，“语言”取有词有语法的系统，音乐和手艺在外面，对反方稍好；“人类的边界”包含“传给别人”，给正方留了积累的战场。整体略偏反。}{略偏反}
\assess{论证义务}{“是”要求基本重合，正方要证明两件事：人超出本能的那一截以语言为条件，语言之外的留不住；反方只要一块稳定、重要的语言外领地。正方更重。}{偏反}
\assess{论据可得性}{成人脑科学和失语研究（Fedorenko 2016、2024；Varley 2005；Monti 2009）站反方；发育研究（Pyers \& Senghas 2009、Spaepen 2011）和传递实验（Morgan 2015）站正方。多条论据双刃：皮拉罕、俄罗斯蓝、轮扁、海伦·凯勒。}{略偏反}
\assess{评委直觉}{没有评委信息，只作一般判断：“言不尽意”“只可意会不可言传”是中文听众的常识，反方先得直觉；维特根斯坦的名言给正方一点开场好感，但被“5.6 说的是我的”拆掉后会反噬。}{偏反}
\assess{赛制顺序}{正方立论先说、结辩最后一发，还能留临场位回应反一；反方先申论（反二）可以先定战场。整体略偏正。}{略偏正}
\end{assessment}
\begin{fields}
\claim[结论]{正方较劣势，劣势主要在论证义务和评委直觉。}
\begin{procard}{劣势方打法}\sub{收窄战场}{只打两处：人超出动物的那一截（数、错误信念），人类的接力（石器、轮扁、性骚扰）。感知、情感、音乐的战场一律承认“那是底子”，不接。}\sub{抢先定义}{一辩稿第一分钟就把“一个人摸到”和“人类走到”分开，并主动承认失语者能思考，再用脚手架接住。}\sub{判准之争}{反复问“稳定地到达，还是一个人摸到”；反方每举一个人的例子，都追一句“他走了以后呢”。}\sub{价值层}{聋孩子：一个人能走多远，取决于有没有人把词递给他。这是正方把“语言的边界”变成“责任”的地方。}\sub{顺序利用}{正一结辩最后一发，回应组按反一可能收的三条路写，临场位只留一句。}\end{procard}
\begin{concard}{优势方要防的}\sub{滑坡}{别说成“语言不重要”“语言束缚人”；反方的立场是语言是最好的工具，只是不是边界。}\sub{脚手架}{失语者被“先有语言”打回，要立刻接“孩子身上两者早就分开”和“语言剥夺者只是延迟”。}\sub{双刃}{轮扁、石器实验都能被正方反用，反方只引对自己有利的那半句时，要准备另一半被念出来。}\end{concard}
\end{fields}

\section{十、口径表}
\prosubsection{10.1 正方口径表}
\begin{canon}{正方}
\canongroup{定义}
\canonrow{语言}{语言，是有词、有语法、用来表达和交流意思的符号系统。口语、文字、手语都是；音乐、绘画、表情叫语言，只是比喻。}{有词有语法的那套东西；嘴上、纸上、手上的语言都算}{语言就是一切表达；音乐也是语言}{汉典；Pyers \& Senghas 2009}
\canonrow{人类的边界}{人能想到、认识到、做到，并且能传给别人的范围的极限。}{人能走多远；人类一代接一代能到的地方}{人类的边界就是人和动物的分界}{汉典“界限”：限度、止境}
\canonrow{是}{两条边界基本重合。语言到不了的地方，人类在重要的方面也到不了。}{两条线大体叠在一起}{是就是影响}{中立定义}
\canongroup{判准}
\canonrow{中立判准}{看语言到不了的地方，人类还能不能稳定地到达。}{语言外面，人是一个人摸到，还是人类走到}{只要证明语言有用就行}{中立判准}
\canonrow{摸到与走到}{一个人摸到，是他的边界；说出来、传下去，人类才走到。}{一个人到了不算人类到了}{个人不算人类}{正方口径}
\canongroup{论点}
\canonrow{论点一\newline 没有词，走不到}{人超出动物的那一截思考，精确的数、别人心里想错了什么，是跟着语言一起到的。语言走到哪，人才走到哪。}{词是思想的脚手架；一辈子缺这个词，就走不到这一步}{没有语言就不能思考；失语的人不会想}{Pyers \& Senghas 2009；维果茨基 1934}
\canonrow{论点二\newline 说不出就接不上}{人类的边界，是一代接一代往前走的边界。说不出的东西能被守住，却接不上力，跟着一个人活，也跟着一个人走。}{说不出的只能守，不能攒；没运到的货随人而去}{说不出就完全传不下}{Morgan 等 2015；Tennie 等 2009}
\canonrow{全场问题}{说不出的东西，怎么让下一个人接着往前走？}{说不出的，下一代从哪里接}{}{一辩稿抛出}
\canongroup{数据}
\canonrow{尼加拉瓜}{第一批 8 人，平均 26.8 岁，4 题平均答对 0.5 题；第二批 10 人平均 3.8 题；作者：二十五年的社会经验也替代不了。}{四道题平均只答对半道；二十五年的社交经验也替代不了}{证明没有语言就没有心智}{Pyers \& Senghas 2009}
\canonrow{家庭手势者}{4 个成年人，有工作、会挣钱；1 到 3 个全对，超过 3 个精确答对不到一半。}{过了三就不准；四个人都认得钱}{家庭手势者不会数数}{Spaepen 等 2011}
\canonrow{石器实验}{184 人；只模仿没有提高，只比划翻一倍，能说话翻四倍；奥杜威石器七十多万年几乎没变，作者推测与传得不准有关。}{一百八十四个人学打石头}{证明古人没有语言所以停滞}{Morgan 等 2015}
\canonrow{传递工具}{Fedorenko 等 2024：语言是传递文化知识的强大工具。}{连反方常引的那篇 Nature 也这么写}{Fedorenko 支持正方}{Fedorenko 等 2024}
\canonrow{聋孩子}{美国聋生里，父母至少一方是聋人的不到 5\%。}{不到二十个里有一个；聋孩子多半生在听人家里}{95\% 的聋童没有语言}{Mitchell \& Karchmer 2004}
\canongroup{案例}
\canonrow{轮扁}{口不能言，有数存焉于其间；臣不能以喻臣之子；古之人与其不可传也死矣。}{七十岁的手艺，儿子学不走}{轮扁证明手艺没法传}{《庄子·天道》}
\canonrow{性骚扰一词}{1975 年，康奈尔的 Carmita Wood 被骚扰多年后辞职，Lin Farley 等人和她开会，定下了这个名字。}{一个名字，把一个人的遭遇变成所有人的知识}{没有这个词之前没人觉得这是错的}{History.com 2018；Fricker 2007}
\canongroup{分工}
\canonrow{二辩申论}{主任务：拆反方立论的主干；次任务：推进论点一的新一层（家庭手势者）。}{}{重复立论}{赛制文件}
\canonrow{三辩申论}{主任务：处理反方全场最强的一点；次任务：推进论点二的新一层（性骚扰一词），为结辩定战场。}{}{和二辩讲同一件事}{赛制文件}
\end{canon}

\consubsection{10.2 反方口径表}
\begin{canon}{反方}
\canongroup{定义}
\canonrow{语言}{语言，是有词、有语法、用来表达和交流意思的符号系统。口语、文字、手语都是；音乐、绘画、表情叫语言，只是比喻。}{有词有语法的那套东西；嘴上、纸上、手上的语言都算}{语言就是嘴上说的话}{汉典；语言学共识}
\canonrow{人类的边界}{人能想到、认识到、做到，并且能传给别人的范围的极限。}{人能走多远；人能想到、做到的地方}{人类的边界就是人能感受到的一切}{汉典“界限”：限度、止境}
\canonrow{是}{两条边界基本重合。语言到不了的地方，人类在重要的方面也到不了。}{两条线大体叠在一起}{只要有一处不重合就不是}{中立定义}
\canongroup{判准}
\canonrow{中立判准}{看语言到不了的地方，人类还能不能稳定地到达。}{语言外面，人是不是一再到得了}{只要有说不清的感受就赢}{中立判准}
\canonrow{工具不是边界}{语言是人类最好的工具；好用的工具让人走得快，不划人能走多远。}{语言是出口，不是围墙}{语言不重要；语言束缚人}{反方口径}
\canongroup{论点}
\canonrow{论点一\newline 想在说前面}{思考不是用语言做的。语言坏了，人照样算数、推理、作曲；语言还没到，人已经想到了。}{先想到，再找词；语言是思想的出口，不是思想的围墙}{语言对思考没用}{Fedorenko 等 2024；Varley 等 2005}
\canonrow{论点二\newline 知道的比说的多}{人最高处的本事，棋感、手上的分寸、医生的直觉，说不清，却真实存在，靠做和示范一代代传下去。说不出，不等于到不了。}{看得出的比叫得出的多；说不出，照样做得出、教得会}{说不出的比说得出的重要}{波兰尼 1966；Morgan 等 2015}
\canonrow{全场问题}{一个中风失语的人，还能不能思考？}{失语的人还能不能想事情}{}{一辩稿抛出}
\canongroup{数据}
\canonrow{失语算术}{三名重度语法失语者，基本运算全部完好，能解带括号的算式。}{说不出一句完整的话，照样会算}{失语者一切正常}{Varley 等 2005}
\canonrow{全面失语}{全面失语者仍能加减、解逻辑题、推断别人的想法、欣赏音乐；一位俄国作曲家严重失语后继续作曲，作品被评为不输病前。}{话说不出来，脑子还在转}{失语者和常人没有区别}{Fedorenko \& Varley 2016}
\canonrow{不是前提}{Fedorenko 等 2024：语言不像是复杂思考的前提。}{语言是说给别人听的工具}{科学已经证明语言和思考无关}{Fedorenko 等 2024}
\canonrow{爱因斯坦}{写出来或说出来的词语，在我的思考机制里似乎不起任何作用；词语要到第二阶段才费力去找。}{他先想到，后找词}{爱因斯坦不用语言思考}{Hadamard 1945 附录二}
\canonrow{俄罗斯蓝}{俄语 26 人、英语 24 人；准确率英语者 96.5\%、俄语者 95.7\%；俄语者快 124 毫秒，言语干扰下消失。}{准确率都在九成六上下；差别只是一百二十多毫秒}{语言对感知毫无影响}{Winawer 等 2007}
\canonrow{比划组}{只比划、不出声的教学组，每下敲出可用石片的概率翻一倍。}{不出声也教得会}{隐瞒说话组翻四倍}{Morgan 等 2015}
\canongroup{案例}
\canonrow{舍巴林}{苏联作曲家舍巴林，中风后严重失语，继续作曲，作品被评为与病前水准相当。}{说不出话的作曲家还在写曲子}{他写出了最伟大的作品}{Luria 等 1965；Fedorenko \& Varley 2016}
\canonrow{认脸}{你能在一千张脸里认出你妈，却说不出她长什么样。}{认得出，说不出}{}{波兰尼 1966}
\canongroup{分工}
\canonrow{二辩申论}{主任务：拆正方立论的主干；次任务：推进论点一的新一层（爱因斯坦的信）。}{}{重复立论}{赛制文件}
\canonrow{三辩申论}{主任务：处理正方全场最强的一点；次任务：推进论点二的新一层（俄罗斯蓝），为结辩定战场。}{}{和二辩讲同一件事}{赛制文件}
\end{canon}

\section{十一、论据出处清单}
\begin{sourcelist}
\source{1}{《逻辑哲学论》5.6：我的语言的界限意味着我的世界的界限}{Wittgenstein, Tractatus Logico-Philosophicus (1921/1922)，Ogden 与 Pears–McGuinness 双译对照本，Klement 编，\url{https://people.umass.edu/klement/tlp/tlp.html}}{Die Grenzen meiner Sprache bedeuten die Grenzen meiner Welt.}{2026-09-30}{已核实}{双方质询、反驳库}{}
\source{2}{《逻辑哲学论》6.522：确有不可说的东西，它显示自身}{Wittgenstein (1922) Tractatus Logico-Philosophicus，Klement 对照本，\url{https://people.umass.edu/klement/tlp/tlp.html}}{Es gibt allerdings Unaussprechliches. Dies zeigt sich, es ist das Mystische.}{2026-09-30}{已核实}{反方反驳库、质询}{}
\source{3}{《逻辑哲学论》5.641：哲学的我不是人}{Wittgenstein (1922) Tractatus，5.641，\url{https://people.umass.edu/klement/tlp/tlp.html}}{Das philosophische Ich ist nicht der Mensch}{2026-09-30}{已核实}{反方反驳库、质询}{}
\source{4}{《逻辑哲学论》序言：界限只能在语言中划出}{Wittgenstein (1922) Tractatus，序言，\url{https://people.umass.edu/klement/tlp/tlp.html}}{Die Grenze wird also nur in der Sprache gezogen werden können}{2026-09-30}{已核实}{正方反驳库}{}
\source{5}{致 Ficker 信（1919）：作品有两部分，没写的那部分才重要}{Klemme, H. F. (1999). Metaphysik und moralische Verbindlichkeit（引 Wittgenstein 1980, S. 96）；SEP “Ludwig Wittgenstein” 英译}{}{}{有把握}{反方反驳库}{原信：Wittgenstein, Briefe an Ludwig von Ficker (1969)，查原文页码}
\source{6}{《庄子·天道》轮扁斫轮}{《庄子·天道》，维基文库本（2026 年访问），\url{https://zh.wikisource.org/wiki/\%E8\%8E\%8A\%E5\%AD\%90/\%E5\%A4\%A9\%E9\%81\%93}}{口不能言，有数存焉于其间。臣不能以喻臣之子}{2026-09-30}{已核实}{正方论点二案例；反方反驳库}{}
\source{7}{《易传·系辞上》：书不尽言，言不尽意；圣人立象以尽意}{《易传·系辞上》，维基文库本（2026 年访问），\url{https://zh.wikisource.org/wiki/\%E6\%98\%93\%E5\%82\%B3/\%E7\%B9\%AB\%E8\%BE\%AD\%E4\%B8\%8A}}{子曰：“书不尽言，言不尽意。”}{2026-09-30}{已核实}{正方反驳库}{}
\source{8}{汉典“语言”“边界”“界限”词条}{汉典 zdic.net（2026 年访问），\url{https://www.zdic.net/}}{人类用来表达思想、进行交际的工具，由语音、词汇和语法构成}{2026-09-30}{已核实}{定义}{}
\source{9}{维果茨基：思想在词里完成}{Vygotsky, L. S. (1934). Thinking and Speech, ch.7（Minick 英译），\url{https://www.marxists.org/archive/vygotsky/works/words/ch07.htm}}{Thought is not expressed but completed in the word.}{2026-09-30}{已核实}{正方论点一学理}{}
\source{10}{尼加拉瓜手语：第一批 8 人（26.8 岁）4 题答对 0.5 题，第二批 10 人 3.8 题；二十五年社会经验替代不了}{Pyers \& Senghas (2009). Psychological Science 20(7), 805–812. doi:10.1111/j.1467-9280.2009.02377.x}{cannot be replaced by even 25 years of social experience}{2026-09-30}{已核实}{正方立论、速查}{}
\source{11}{家庭手势者 4 人（20–29 岁），1 到 3 个全对，超过 3 个精确答对不到一半；他们有工作、会挣钱}{Spaepen et al. (2011). PNAS 108(8), 3163–3168. doi:10.1073/pnas.1015975108}{do not spontaneously develop representations of large exact numerosities}{2026-09-30}{已核实}{正二申论}{}
\source{12}{石器实验：184 人；模仿没提高、比划翻一倍、说话翻四倍；奥杜威停滞七十多万年（700,000 年）}{Morgan et al. (2015). Nature Communications 6, 6029. doi:10.1038/ncomms7029}{gestural teaching doubled and verbal teaching quadrupled this probability}{2026-09-30}{已核实}{正方立论、反方论点二}{}
\source{13}{棘轮效应}{Tennie, Call \& Tomasello (2009). Phil. Trans. R. Soc. B 364, 2405–2415. doi:10.1098/rstb.2009.0052}{accumulates modifications over time (what we call the 'ratchet effect')}{2026-09-30}{已核实}{正方论点二学理}{}
\source{14}{语言不是复杂思考的前提；语言是传递文化知识的强大工具}{Fedorenko, Piantadosi \& Gibson (2024). Nature 630, 575–586. doi:10.1038/s41586-024-07522-w}{language does not appear to be a prerequisite for complex thought}{2026-09-30}{已核实}{反方论点一；正方论点二}{}
\source{15}{全面失语者仍能加减、解逻辑题、推断他人想法、欣赏音乐；俄国作曲家失语后继续作曲}{Fedorenko \& Varley (2016). Annals of the NYAS 1369, 132–153. doi:10.1111/nyas.13046}{able to add and subtract, solve logic problems, think about another person's thoughts}{2026-09-30}{已核实}{反方立论、速查}{}
\source{16}{三名重度语法失语者运算完好}{Varley et al. (2005). Agrammatic but numerate. PNAS 102(9), 3519–3524. doi:10.1073/pnas.0407470102}{all basic computational procedures were intact across patients}{2026-09-30}{已核实}{反方立论}{}
\source{17}{演绎推理不在语言区}{Monti, Parsons \& Osherson (2009). The boundaries of language and thought in deductive inference. PNAS 106, 12554–12559. doi:10.1073/pnas.0902422106}{logical inference is not embedded in natural language}{2026-09-30}{已核实}{反方备用}{}
\source{18}{舍巴林失语后作曲（书目）}{Luria, Tsvetkova \& Futer (1965). Aphasia in a composer (V. G. Shebalin). J. Neurol. Sci. 2, 288–292. doi:10.1016/0022-510x(65)90113-9}{}{}{有把握}{反方案例（内容以第 15 条为准）}{查原文正文，确认作曲时间与作品}
\source{19}{舍巴林失语后完成第五交响曲，肖斯塔科维奇称赞}{只见维基百科等二手来源}{}{}{待核实}{备用（未进稿件）}{查 Luria 1965 原文或舍巴林传记}
\source{20}{爱因斯坦：词语在思考机制中似乎不起作用，第二阶段才费力去找}{Hadamard (1945). The Psychology of Invention in the Mathematical Field. Princeton UP（Dover 1954），附录二 pp.142–143，\url{https://worrydream.com/refs/Hadamard_1945_-_The_psychology_of_invention_in_the_mathematical_field.pdf}}{do not seem to play any role in my mechanism of thought}{2026-09-30}{已核实}{反二申论}{}
\source{21}{俄罗斯蓝：俄语 26 人、英语 24 人；准确率 96.5\% 与 95.7\%；类别优势 124 毫秒}{Winawer et al. (2007). PNAS 104(19), 7780–7785. doi:10.1073/pnas.0701644104}{accuracy was high (96.5 ± 2.1\% and 95.7 ± 3.2\%}{2026-09-30}{已核实}{反三申论}{}
\source{22}{五个月婴儿会算小数目加减}{Wynn (1992). Addition and subtraction by human infants. Nature 358, 749–750. doi:10.1038/358749a0}{5-month-old infants can calculate the results of simple arithmetical operations}{2026-09-30}{已核实}{反方备用}{}
\source{23}{皮拉罕：没有表达精确数量的办法；实物在场可精确配对，涉及记忆不准}{Frank, Everett, Fedorenko \& Gibson (2008). Cognition 108, 819–824. doi:10.1016/j.cognition.2008.04.007}{no linguistic method whatsoever for expressing exact quantity, not even "one."}{2026-09-30}{已核实}{反驳库}{}
\source{24}{皮拉罕：超过三的数量表现很差}{Gordon (2004). Science 306, 496–499. doi:10.1126/science.1094492}{Performance with quantities greater than three was remarkably poor}{2026-09-30}{已核实}{正方备用}{}
\source{25}{默会知识：我们知道的比能说出来的多；认脸}{Polanyi (1966). The Tacit Dimension. University of Chicago Press, ch.1}{}{}{有把握}{反方论点二学理、案例}{查原书第 4 页原句}
\source{26}{海伦·凯勒：老师来之前活在“无世界”里；那时会生气、满足、渴望}{Keller (1908). The World I Live In. Century Co., ch. XI “Before the Soul Dawn”，\url{https://archive.org/details/worldilivein00kellgoog}}{I lived in a world that was a no-world.}{2026-09-30}{已核实}{双方反驳库}{}
\source{27}{沃尔夫：霍皮语没有直接指时间的词和语法形式}{Whorf (1956). Language, Thought, and Reality. MIT Press，\url{https://archive.org/details/languagethoughtr00whor}}{contain no words, grammatical forms, constructions or expressions that refer directly to what we call "time,"}{2026-09-30}{已核实}{反方备用}{}
\source{28}{Malotki 1983 用 600 多页实例反驳沃尔夫的霍皮时间说}{Malotki (1983). Hopi Time. Mouton；Comrie 书评，见 Malotki 主页 \url{https://ac.nau.edu/~malotki/html/hopi_time_reviews.html}}{Malotki's monograph is a detailed refutation of the key points that Whorf makes}{2026-09-30}{已核实}{反方备用}{}
\source{29}{诠释不正义；本人也可能完全清楚发生了什么}{Fricker (2013). Epistemic Justice as a Condition of Political Freedom? Synthese 190(7), 1317–1332，作者主页 PDF \url{https://www.mirandafricker.com}}{Hermeneutical injustice occurs at a stage prior to communicative activity}{2026-09-30}{已核实}{双方反驳库}{}
\source{30}{诠释不正义与性骚扰案例}{Fricker (2007). Epistemic Injustice. Oxford UP, ch.7}{}{}{有把握}{正方三辩申论学理}{查第 7 章 Carmita Wood 一段}
\source{31}{性骚扰一词 1975 年在康奈尔定名；Farley：我们需要给这个现象一个名字}{Blakemore, E. (2018). Until 1975, “Sexual Harassment” Was the Menace With No Name. History.com, \url{https://www.history.com/articles/until-1975-sexual-harassment-was-the-menace-with-no-name}}{I thought that we needed to have a name for what this phenomenon was.}{2026-09-30}{已核实}{正三申论；反方反驳库}{}
\source{32}{Searle：一切人类制度实在由地位功能宣告创造和维持}{Searle (2010). Making the Social World. Oxford UP, p.201；经 Tsohatzidis 2010 NDPR 书评引}{}{}{有把握}{正方备用}{查原书第 201 页}
\source{33}{美国聋生里父母至少一方是聋人的不到 5\%}{Mitchell \& Karchmer (2004). Chasing the Mythical Ten Percent. Sign Language Studies 4(2), 138–163，ERIC EJ747626 \url{https://eric.ed.gov/?id=EJ747626}}{less than five percent of deaf and hard of hearing students}{2026-09-30}{已核实}{正方结辩价值升华}{}
\source{34}{推迟接触手语的聋童，心理理论脑区像更小的孩子}{Richardson et al. (2020). Nature Communications 11, 3246. doi:10.1038/s41467-020-17004-y}{Early linguistic experience may facilitate ToM development}{2026-09-30}{已核实}{正方备用}{}
\source{35}{边跟读边找方向的成人表现得像幼儿和老鼠}{Hermer-Vazquez, Spelke \& Katsnelson (1999). Cognitive Psychology 39, 3–36. doi:10.1006/cogp.1998.0713}{reoriented themselves like children and adult rats}{2026-09-30}{已核实}{不上场（见第 36 条）}{}
\source{36}{上一条未能复现，语言不是必需}{Ratliff \& Newcombe (2008). Cognitive Psychology 56, 142–163. doi:10.1016/j.cogpsych.2007.06.002}{language is not necessary for successful use of features in reorientation}{2026-09-30}{已核实}{反方备用}{}
\source{37}{《哲学研究》§23：说语言是一种活动、一种生活形式的一部分}{Wittgenstein (1953). Philosophical Investigations §23；经 SEP “Ludwig Wittgenstein” 引}{}{}{有把握}{备用}{查 Anscombe 译本 §23}
\source{38}{达尼人只有两个基本颜色词，颜色记忆与英语者相近}{Heider, E. R. (1972). Universals in color naming and memory. J. Exp. Psychol. 93, 10–20（只见书目，摘要未开）}{}{}{待核实}{备用（未进稿件）}{查原文结果部分；注意 Roberson 等 2000 的反例}
\source{39}{野外恒河猴也能察觉小数目的加减变化}{Hauser, MacNeilage \& Ware (1996). Numerical representations in primates. PNAS 93, 1514–1517. doi:10.1073/pnas.93.4.1514}{rhesus monkeys detect additive and subtractive changes in the number of objects}{2026-09-30}{已核实}{正方防守与对辩“猴子身上也测到过”}{}
\end{sourcelist}

\section{十二、备赛建议}
\begin{fields}
\begin{timeline}\when{第 1 天}{全队对齐定义、中立判准和“摸到与走到”（正方）/“工具不是边界”（反方）两句口径。}\when{第 2–3 天}{一辩背立论；二辩、三辩按分工练申论套件，每人念三种组装。}\when{第 4–5 天}{质询与被质询对练：90 秒双边，计打断次数；对辩三场各练两遍。}\when{第 6 天}{完整模辩一场，按记录卡组装申论与结辩，赛后更新口径表。}\end{timeline}
\field{模辩重点检验}{\sub{一}{正方：被问“失语的人还能不能思考”时，三个人给出同一个答案。}\sub{二}{反方：被念轮扁后半句、石器“说话翻四倍”时，三个人给出同一个答案。}\sub{三}{双方：申论里的每个数字，下一分钟被质询时能报出作者和年份。}}
\begin{checklist}\item 中立定义三句（语言、人类的边界、是）\item 中立判准一句和双方义务\item 我方两个论点标签与一句话\item 全场问题\item 对方最可能打的两处和我方口径\item 我方每个数字的作者与年份\end{checklist}
\end{fields}
\end{document}
